\begin{frame}{Week 13의 전제가 실차 규모에서 만나는 곳}
  \begin{itemize}
    \item Week 13의 전제: ``\textbf{시뮬레이터 안에서 먼저
      익히고, 이후에 실물로 옮긴다}''
    \item Week 13에서 \kb{현실 격차(reality gap)}와
      \kb{도메인 무작위화}를 \textbf{작은 로봇}의 맥락에서
      다룸
  \end{itemize}
  \vskip0.6em
  \begin{block}{자율주행에서의 격차}
    ``실물''이 \textbf{함부로 실험할 수 없는 도로}이기 때문에,
    Offboard 시뮬레이션은 이 책의 이전 어떤 예시보다
    \textbf{``실물에 닮아 있어야''} 평가의 의미가 있다.
    $\Rightarrow$ 현실 격차가 \textbf{전문의 핵심 문제}가 됨.
  \end{block}
\end{frame}
